% Crear pdf con:
% pdflatex -synctex=1 -interaction=nonstopmode %.tex
%"The PDF file may contain up to 25 pages of reference material, single-sided, letter or A4 size, with text and illustrations readable by a person with correctable eyesight without magnification from a distance of 1/2 meter."
%%% hidelinks quita los recuadros rojos de los links.
\documentclass[10pt,landscape,twocolumn,a4paper,notitlepage, hidelinks]{article} 
\usepackage{hyperref}
\usepackage[spanish, activeacute]{babel}
\usepackage[utf8]{inputenc}
\usepackage{fancyhdr}
\usepackage{lastpage}
\usepackage{listings}
\usepackage{amssymb}
\usepackage[usenames,dvipsnames]{color}
\usepackage{graphicx}
\usepackage{wrapfig}
\usepackage{amsmath}
\usepackage{makeidx}

%%% Margenes
\setlength{\columnsep}{0.25in}    % default=10pt
\setlength{\columnseprule}{0.5pt}    % default=0pt (no line)
 
\addtolength{\textheight}{2.35in}
\addtolength{\topmargin}{-0.9in}     % ~ -0.5 del incremento anterior
 
\addtolength{\textwidth}{1.1in}
\addtolength{\oddsidemargin}{-0.55in} % -0.5 del incremento anterior
 
\setlength{\headsep}{0.08in}
\setlength{\parskip}{0in}
\setlength{\headheight}{15pt}
\setlength{\parindent}{0mm}
 
%%% Encabezado y pie de pagina
\pagestyle{fancy}
\fancyhead[LO]{\textbf{\title}}
\fancyhead[C]{\leftmark\ -\ \rightmark}
\fancyhead[RO]{P\'agina \thepage\ de \pageref{LastPage}}
\renewcommand{\headrulewidth}{0.4pt}
\fancyfoot{}
\definecolor{darkblue}{rgb}{0,0,0.4}
%%% Configuracion de Listings
\lstloadlanguages{C++}
\lstnewenvironment{code}
	{%\lstset{	numbers=none, frame=lines, basicstyle=\small\ttfamily, }%
	 \csname lst@SetFirstLabel\endcsname}
	{\csname lst@SaveFirstLabel\endcsname}
\lstset{% general command to set parameter(s)
	language=C++, basicstyle=\small\ttfamily, keywordstyle=\slshape,
	emph=[1]{usa}, emphstyle={[1]\sffamily\bfseries},
	morekeywords={tint,forn,forsn,node,alt,tipo},
	basewidth={0.47em,0.40em},
	columns=fixed, fontadjust, resetmargins, xrightmargin=5pt, xleftmargin=15pt,
	flexiblecolumns=false, tabsize=2, breaklines,	breakatwhitespace=false, extendedchars=true,
	numbers=left, numberstyle=\tiny, stepnumber=1, numbersep=9pt,
	frame=l, framesep=3pt,
    basicstyle=\ttfamily,
    keywordstyle=\color{darkblue}\ttfamily,
    stringstyle=\color{magenta}\ttfamily,
    commentstyle=\color{RedOrange}\ttfamily,
    morecomment=[l][\color{OliveGreen}]{\#}
}

\lstdefinestyle{C++}{
	language=C++, basicstyle=\small\ttfamily, keywordstyle=\slshape,
	emph=[1]{usa,tipo2}, emphstyle={[1]\sffamily\bfseries},
	morekeywords={tint,forn,forsn,tipo,alt,node},
	basewidth={0.47em,0.40em},
	columns=fixed, fontadjust, resetmargins, xrightmargin=5pt, xleftmargin=15pt,
	flexiblecolumns=false, tabsize=2, breaklines,	breakatwhitespace=false, extendedchars=true,
	numbers=left, numberstyle=\tiny, stepnumber=1, numbersep=9pt,
	frame=l, framesep=3pt,
    basicstyle=\ttfamily,
    keywordstyle=\color{darkblue}\ttfamily,
    stringstyle=\color{magenta}\ttfamily,
    commentstyle=\color{RedOrange}\ttfamily,
    morecomment=[l][\color{OliveGreen}]{\#}
}
 
%%% Macros
\def\nbtitle#1{\begin{Large}\begin{center}\textbf{#1}\end{center}\end{Large}}
\def\nbsection#1{\section{#1}}
\def\nbsubsection#1{\subsection{#1}}
\def\nbcoment#1{\begin{small}\textbf{#1}\end{small}}
\newcommand{\comb}[2]{\left( \begin{array}{c} #1 \\ #2 \end{array}\right)}
\def\complexity#1{\texorpdfstring{$\mathcal{O}(#1)$}{O(#1)}}
 \newcommand\cppfile[2][]{
\lstinputlisting[style=C++,linerange={#1}]{#2}
}

%%% Configuración de paths

\makeatletter
\def\input@path{{../}}
\makeatother

%%% Arregla espacios

\frenchspacing

\begin{document}
\def\title{Universidad Nacional del Sur - Mastropiero Limit Exceeded}

%\begin{center}{\includegraphics[width=4cm]{../others/mastro}}\end{center}

%\tableofcontents\newpage
 
%\section{Referencia}%%%%%%%%%%%%%%%%%%REFERENCIA%%%%%%%%%%%%%%%%%%
%\begin{tabular}{|l|l|p{5.4cm}|} \hline
\textbf{Algorítmo} & \textbf{Parámetros} &  \textbf{Función} \\  \hline
%swap & e1, e2 &  da vuelta e1,e2 & $1$\\\hline
sort, stable\_sort & f, l &  ordena el intervalo \\  \hline
%is\_sorted & f, l &  \textit{bool} si esta ordenado \\  \hline
nth\_element & f, nth, l & \textit{void} ordena el n-esimo, y \\ && particiona el resto \\  \hline
fill, fill\_n & f, l / n, elem & \textit{void} llena [f, l) o [f, \\ && f+n) con elem \\  \hline
lower\_bound, upper\_bound & f, l, elem & \textit{it} al primer / ultimo donde se \\ && puede insertar elem para que\\ && quede ordenada \\  \hline
binary\_search & f, l, elem & \textit{bool} esta elem en [f, l) \\  \hline
copy & f, l, resul & hace resul+$i$=f+$i$ $\forall i$ \\  \hline
find, find\_if, find\_first\_of & f, l, elem & \textit{it} encuentra i $\in$[f,l) tq. i$=$elem, \\ & / pred / f2, l2 & pred(i), i$\in$[f2,l2)\\\hline
count, count\_if & f, l, elem/pred & cuenta elem, pred(i)\\\hline
search & f, l, f2, l2 & busca [f2,l2) $\in$ [f,l)\\\hline
replace, replace\_if & f, l, old & cambia old / pred(i) por new \\ & / pred, new &\\\hline
reverse & f, l & da vuelta\\\hline
partition, stable\_partition & f, l, pred & pred(i) ad, !pred(i) atras\\\hline
%min, max & e1, e2 & men / may & $1$\\\hline
min\_element, max\_element & f, l, [comp] & \textit{it} min, max de [f,l]\\\hline
lexicographical\_compare & f1,l1,f2,l2 & \textit{bool} con [f1,l1]<[f2,l2]\\\hline
next/prev\_permutation & f,l & deja en [f,l) la perm sig, ant\\\hline
set\_intersection, & f1, l1, f2, l2, res & [res, $\ldots$) la op. de conj\\
set\_difference, set\_union, & & \\
set\_symmetric\_difference, & &\\\hline
push\_heap, pop\_heap, & f, l, e / e / & mete/saca e en heap [f,l), \\
make\_heap & & hace un heap de [f,l)\\\hline
is\_heap & f,l & \textit{bool} es [f,l) un heap\\\hline
accumulate & f,l,i,[op] & \textit{T} $=$ $\sum$/oper de [f,l)\\\hline
inner\_product & f1, l1, f2, i & \textit{T} $=$ i $+$ [f1, l1) . [f2, $\ldots$ )\\\hline
partial\_sum & f, l, r, [op] & r+i = $\sum$/oper de [f,f+i] $\forall i \in$[f,l)\\\hline
%power & e, i, op & \textit{T} = $e^{n}$\\\hline
\_\_builtin\_ffs& unsigned int & Pos. del primer 1 desde la derecha\\\hline
\_\_builtin\_clz & unsigned int & Cant. de ceros desde la izquierda.\\\hline
\_\_builtin\_ctz & unsigned int & Cant. de ceros desde la derecha.\\\hline
\_\_builtin\_popcount & unsigned int & Cant. de 1’s en x.\\\hline
\_\_builtin\_parity & unsigned int & 1 si x es par, 0 si es impar.\\\hline
\_\_builtin\_XXXXXXll & unsigned ll & = pero para long long's.\\\hline
\end{tabular}\newpage


\section{Estructuras}%%%%%%%%%%%%%%%%%%ESTRUCTURAS%%%%%%%%%%%%%%%%%%
%\subsection{Sparse Table}
%\cppfile{../structures/sparse.table.cpp}
%\subsection{Segment Tree}
%\cppfile{structures/segment.tree.cpp}
%\subsection{Segment Tree (Iterative)}
%\cppfile{structures/segment.tree.iterative.cpp}
\subsection{Segment Tree (Lazy)}
\cppfile{structures/segment.tree.lazy.cpp}
%\subsection{Segment Tree (Persistent)}
%\cppfile{structures/segment.tree.persistent.cpp}
%\subsection{Sliding Window RMQ}
%\cppfile{structures/sliding.window.rmq.cpp}
%\subsection{Fenwick Tree}
%\cppfile{structures/fenwick.cpp}
%
%\subsection{Disjoint Intervals}
%\cppfile{structures/disjoint.intervals.cpp}
%\subsection{Segment Tree (2D)}
%\cppfile{structures/segment.tree.2d.cpp}
%\subsection{Big Int}
%\cppfile{structures/bigint.cpp}
%\subsection{Modnum}
%\cppfile{structures/mod.cpp}

\subsection{Treap}
%\textbf{Definición:} estructura de datos que combina los conceptos de \emph{binary search tree} (para las claves)
%y \emph{heap} (para las prioridades), y asigna las prioridades de forma aleatoria para asegurar una altura de 
%$O(\log{}n)$ en promedio.
%
%\textbf{Operaciones básicas:}
%\begin{itemize}
%	\item $split(T, X)$: separa al árbol $T$ en 2 subárboles $T_L$ y $T_R$ tales que $T_L$ contiene a todos los 
%	elementos con claves menores a $X$ y $T_R$ a los demás.
%	\item $merge(T_1, T_2)$: combina dos subárboles $T_1$ y $T_2$ y retorna un nuevo árbol, 
%	asume que las claves en $T_1$ son menores que las claves en $T_2$.
%\end{itemize}
%
%\textbf{Operaciones avanzadas:}
%\begin{itemize}
%	\item $insert(T, X)$: inserta una nueva clave al árbol. Resulta trivial de implementar a partir de las anteriores:
%	$(T_1, T_2) = split(T, X)$ y $T_3 = merge(merge(T_1, X), T_2)$.
%\end{itemize}
%\subsubsection{Treap set}
%\cppfile{structures/treap.set.cpp}
%\subsubsection{Treap array}
%\textbf{Explicación treap implícito:} permite insertar, borrar, hacer queries y updates (incluyendo reverse) en rangos
%en un arreglo. La idea es usar a los índices como claves, pero en vez de almacenarlos (sería difícil actualizar en ese caso),
%aprovechamos que la clave de un nodo es la cantidad de elementos menores a ese nodo (cuidado, no son solo los del
%subárbol izquierdo).
\cppfile{structures/treap.array.cpp}

\subsection{Convex Hull Trick}
\cppfile{structures/chull.trick.cpp}
\subsection{Convex Hull Trick (Dynamic)}
\cppfile{structures/chull.trick.dyn.cpp}
%\subsection{Li-Chao Tree}
%\cppfile{structures/li.chao.tree.cpp}
%\subsection{Gain-Cost Set}
%\cppfile[1-24]{structures/gain.cost.set.cpp}
%\subsection{Set con índices}
%\cppfile{structures/order.statistics.tree.cpp}


%\section{Algoritmos varios}%%%%%%%%%%%%%%%%%%ALGORITMOS%%%%%%%%%%%%%%%%%%%%%%%%%%
%\subsection{Longest Increasing Subsequence}
%\cppfile[1-31]{misc/lis.cpp}
%\subsection{Alpha-Beta prunning}
%\cppfile{misc/alphabeta.cpp}
%\subsection{Mo's algorithm}
%\cppfile{misc/mosalgorithm.cpp}
%\subsection{Parallel binary search}
%\textbf{Descripción:} permite reutilizar información cuando se necesitan
%realizar múltiples búsquedas binarias sobre la misma información.
%
%\textbf{Explicación algoritmo:} imaginarse un árbol binario de rangos de
%búsqueda binaria $(lo, hi]$ y queries asignadas a cada nodo, que implican que
%esa query está en ese rango de la búsqueda binaria. El algoritmo aprovecha
%que para cada nivel del árbol las queries están ordenadas, y se puede procesar
%la información hasta el $mid$ de cada query en orden, resultando en un tiempo de
%$O(N + Q_{nivel})$ por nivel (más un $log$ extra por ordenar). 
%
%\textbf{Observación: } se puede implementar de forma recursiva, dependiendo del
%problema. Esto puede mejorar la complejidad ya que se evita el ordenamiento.
%\cppfile{misc/parallel_binary_search.cpp}

\section{Strings}%%%%%%%%%%%%%%%%%%STRINGS%%%%%%%%%%%%%%%%%%%%%%%%%%
\subsection{Hash}
\cppfile{strings/hash.cpp}
\subsection{Manacher}
\textbf{Definición:} permite calcular todas las substrings
de una string $s$ que son palíndromos.
Para ello, mantiene un arreglo $odd$ tal que $odd[i]$ 
almacena la longitud del palíndromo impar maximal con centro en $i$. 
Análogamente mantiene un arreglo $even$ tal que $even[i]$ guarda la 
longitud del palíndromo par maximal con centro derecho en $i$. \\
\textbf{Explicación del algoritmo:} muy similar al algoritmo para calcular
la \emph{función Z}, mantiene el palíndromo que termina más a la derecha
entre todos los palíndromos ya detectados con rango $[l, r]$. Utiliza la 
información ya calculada si $i$ está dentro de $[l, r]$, y luego corre
el algoritmo trivial. Cada vez que se corre el algoritmo trivial, r se 
incrementa en 1 y r jamás decrece.\\
\cppfile{strings/manacher.cpp}
\subsection{KMP}
\cppfile{strings/kmp.cpp}
\subsection{Trie}
\cppfile{strings/trie.cpp}
\subsection{Suffix Array (corto, nlog2n)}
\cppfile{strings/suffix.array.short.cpp}
%\subsection{Suffix Array (largo, nlogn)}
%\cppfile[1-38]{strings/suffix.array.cpp}
\subsection{String Matching With Suffix Array}
\cppfile[38-59]{strings/suffix.array.cpp}
\subsection{LCP (Longest Common Prefix)}
\cppfile[61-75]{strings/suffix.array.cpp}
\subsection{Aho-Corasick}
%\textbf{Definición}
%El automáta Aho-Corasick es un autómata $A$ que reconoce un conjunto de cadenas $S$.
%
%\textbf{Conceptos importantes}
%\begin{itemize}
%	\item Cada nodo del autómata se asocia con (al menos) un prefijo de una cadena en $S$.
%	\item Un \emph{suffix link} para un vértice $p$ es un arco que apunta al sufijo propio
%	más largo de la cadena correspondiente al vértice $p$.
%	\item Estando en un estado $p$ que corresponde a una palabra $t$, se pueden definir 
%	arcos de dos tipos:
%	\begin{itemize}
%		\item Transiciones tipo trie: dado un caracter $c$ tal que $t + c$ pertenece al autómata,
%		el arco apunta a $t + c$.
%		\item Transiciones tipo suffix link: dado un caracter $c$ tal que $t + c$ no pertenece al
%		autómata, el arco apunta al máximo sufijo propio de $t + c$ que pertenece al árbol.
%	\end{itemize}
%	\item Implementación:
%	\begin{itemize}
%		\item Cada nodo mantiene:
%		\begin{itemize} 
%			\item Un indicador de la cantidad de cadenas que terminan en ese nodo: \emph{terminal}.
%			\item El padre \emph{p} y el caracter desde el que transicionó \emph{pch}.
%			\item Las transiciones tipo trie en \emph{next}.
%			\item El suffix link en \emph{link}.
%			\item Todas las transiciones (tipo trie y tipo suffix link) en \emph{go}.
%		\end{itemize}
%		\item El algoritmo se divide en:
%		\begin{itemize}
%			\item \emph{add\_string}: agrega una cadena $s$ al autómata.
%			\item \emph{go}: calcula el nodo destino de la transición $(v, ch)$.
%			\item \emph{get\_link}: calcula el suffix link de la cadena correspondiente al nodo $v$.
%		\end{itemize}
%	\end{itemize}
%\end{itemize}
%
%\textbf{Problemas clásicos}
%\begin{itemize}
%	\item Encontrar todas las cadenas de un conjunto en un texto: mantener \emph{exit link}
%	(nodo terminal más cercano alcanzable mediante suffix links), recorrer autómata con el 
%	texto como entrada y transicionar por exit links para encontrar matches.
%	\item Cadena lexicográficamente mínima de longitud $len$ que no matchea ninguna 
%	cadena de un conjunto $S$: DFS sobre autómata para encontrar camino de longitud $L$
%	evitando entrar en nodos terminales.
%	\item Mínima cadena que contiene todas las cadenas de un conjunto $S$: BFS sobre 
%	autómata manteniendo máscara de cadenas matcheadas (y máscara de terminales,
%	incluyendo alcanzables por suffix link, en cada nodo). 
%	\textbf{Recordatorio importante: } un nodo solo mantiene los matches para la cadena
%	completa. Para mantener todos los matches (incluyendo sufijos) estando en un nodo 
%	$v$, hay que usar la información que propagan los suffix links.
%	\item Cadena lexicogrficamente mínima de longitud $len$ que contiene $k$ cadenas
%	de un conjunto $S$: DFS sobre grafo $(v, len, cnt)$.
%\end{itemize}

\cppfile{strings/aho-corasick.cpp}
\subsection{Suffix Automaton}
%\textbf{Definición}
%Un suffix automaton $A$ es un autómata minimal que reconoce los sufijos de una 
%cadena $s$.
%
%\textbf{Conceptos importantes}
%\begin{itemize}
%	\item $A$ \emph{reconoce} a una cadena $s$ si comenzando desde el nodo inicial llegamos a un 
%	terminal.
%
%	\item Dada una subcadena $t$ de $s$, definimos $endpos(t)$ como el conjunto de las posiciones
%	en $s$ en las que terminan las apariciones de $t$.
%	
%    \item Dos subcadenas $u$ y $v$ de $s$ son \emph{equivalentes} si recorrer el autómata con $u$ 
%    y con $v$ nos lleva al mismo nodo. Esto es equivalente a $endpos(u) = endpos(v)$. Los nodos del 
%    automáta se corresponden al conjunto de cadenas de las \emph{clases de equivalencia} bajo la
%    relación anterior.
%
%    \item Las cadenas en una clase de equivalencia son sufijos de la cadena de mayor tamaño de la 
%    clase, y forman un intervalo contiguo de tamaños. El \emph{suffix link} nos lleva al primer sufijo 
%    que no pertenece a esta clase.
%    
%    \item \emph{Suffix tree} implícito (de $s'$): el suffix link saliente de un nodo nos lleva al padre 
%    en el suffix tree de $s'$ y los suffix links entrantes de un nodo provienen de los hijos del suffix 
%    tree de $s'$.
%    
%\end{itemize}
%
%\textbf{Algoritmo para construcción}
%\begin{itemize}
%	\item Agregamos un caracter a la vez. Sea $c$ el caracter a agregar.
%	\item Sea $last$ el estado que corresponde a la cadena entera antes de agregar a $c$.
%	\item Creamos un nuevo estado $cur$, que corresponde a la cadena luego de agregar a $c$.
%	\item Agregamos transiciones a través de $c$ a los sufijos de la cadena (recorriendo suffix 
%	links a partir de $last$), hasta encontrar un estado de un sufijo que ya tenga una transición
%	con $c$.
%	\begin{itemize}
%		\item Si no encontramos un estado, el suffix link de $cur$ es $t_0$.
%		\item Si la transición lleva a un estado $q$ que representa una cadena con un solo caracter más,
%		el suffix link de $cur$ es $q$.
%		\item Si no, es necesario dividir el estado $q$, ya que debemos usarlo como suffix link pero tiene
%		sufijos extra. Después de esto hace falta actualizar los estados que tenían transiciones a $q$.
%	\end{itemize}
%\end{itemize}
%
%\textbf{Problemas clásicos}
%\begin{itemize}
%    \item Determinar si $w$ es subcadena de $s$: simplemente correr el autómata.
%    \item Determinar si $w$ es sufijo de $s$: correr el autómata y ver si caemos en un terminal.
%    \item Contar cantidad de subcadenas distintas de $s$: esto es igual a la cantidad de
%    caminos en el autómata y se calcula mediante una DP.
%    \item Contar cantidad de apariciones de $w$ en $s$: correr autómata con $w$. Llamemos $u$
%    al nodo en el que terminamos, la cantidad de apariciones es la cantidad de caminos 
%    en $A$ que comienzan en $u$ y llegan a un terminal.
%    \item Encontrar dónde aparece $w$ por primera vez en $s$: correr autómata con $w$. Llamemos $u$
%    al nodo en el que terminamos, esto equivale a calcular el camino más largo del autómata 
%    a partir del nodo $u$. Otra solución: mantener $firstpos(v)$, la primera aparición de una subcadena
%    en la cadena (se actualiza cuando se crea un nuevo nodo y cuando se clonan nodos).
%    \item Encontrar las posiciones de todas las apariciones de $w$ en $s$: encontrar el nodo $u$ que 
%    corresponde a $w$, armar el \emph{suffix tree} (mantener los suffix links invertidos), encontrar 
%    todos los nodos en el subárbol con raíz en $u$, cada nodo corresponde a por lo menos una aparición 
%    y cada aparición corresponde a un nodo y su clon (utilizar $firstpos(v)$ para saber la posición, saltear 
%    nodos clonados; o bien agregar un \$ al comienzo de la cadena y encontrar todas las hojas, la posición es
%    la longitud). 
%    \item Subcadena común más larga de un conjunto de cadenas: dadas $k$ cadenas $S_i$, elegimos
%    $k$ separadores distintos entre sí $D_i$, formamos $T = S_1 + D_1 + \dots + S_k + D_k$ y construimos
%    el autómata de esa cadena. Saber si una subcadena pertenece a una cadena $S_i$ en particular corresponde 
%    a verificar que existe un camino a $D_i$ sin pasar por los demás separadores. Si calculamos para cada nodo
%    a qué separadores puede llegar, la respuesta es la máxima de las cadenas más largas de las clases
%    correspondientes a estados $v$ que puede llegar a todos los separadores.
%\end{itemize}
\cppfile{strings/suffix.automaton.cpp}
\subsection{Z Function}
%\textbf{Definición}
%La función Z para una string $s$ de longitud $n$ es un arreglo $a$ de la misma longitud
%tal que $a[i]$ es la \emph{máxima cantidad de caracteres} comenzando desde la posición $i$
%que coinciden con los primeros caracteres de $s$. Es decir, es el \emph{máximo prefijo común}.
%
%\textbf{Observación}
%$z[0]$ no está bien definido, pero se asume igual a $0$.
%
%\textbf{Algoritmo}
%La idea es mantener el máximo \emph{match} (es decir, el segmento $[l, r]$ con
%máximo $r$ tal que se sabe que $s[0..r-l] = s[l..r]$). 
%
%Siendo $i$ el índice actual (del que queremos calcular la función Z), el algoritmo 
%se divide en dos casos:
%\begin{itemize}
%    \item $i > r$: la posición está fuera de lo que hemos procesado. Se corre el
%    \emph{algoritmo trivial}.
%    \item $i <= r$: la posición está dentro del \emph{match actual}, por lo que
%    se puede utilizar como aproximación inicial $z[i] = min(r - i + 1, z[i-l])$,
%    y luego correr el \emph{algoritmo trivial}.
%\end{itemize}
%
%\textbf{Problemas clásicos}
%\begin{itemize}
%    \item Buscar una subcadena: concatenamos $p$ con $t$ (utilizando un separador).
%    Hay una aparición si la función Z matcheó tantos caracteres como la longitud
%    de $p$.
%\end{itemize}
\cppfile{strings/zfunction.cpp}
%\subsection{Palindrome}
%\cppfile{strings/palindrome.cpp}

\section{Geometría}%%%%%%%%%%%%%%%%%%geometria%%%%%%%%%%%%%%%%%%%%%%
\subsection{Epsilon}
\cppfile{geometry/eps.cpp}
\subsection{Point}
\cppfile{geometry/pto.cpp}
\subsection{Orden radial de puntos}
\cppfile{geometry/radial.order.cpp}
%\subsection{Line}
%\cppfile{geometry/line.cpp}
%\subsection{Segment}
%\cppfile{geometry/segm.cpp}
%\subsection{Rectangle}
%\cppfile{geometry/rect.cpp}
%\subsection{Polygon Area}
%\cppfile{geometry/area.cpp}
%\subsection{Circle}
%\cppfile{geometry/circle.cpp}
\subsection{Point in Poly}
\cppfile{geometry/point.in.poly.cpp}
%\subsection{Point in Convex Poly log(n)}
%\cppfile{geometry/point.in.convex.poly.cpp}
%\subsection{Convex Check CHECK}
%\cppfile{geometry/convex.check.cpp}
\subsection{Convex Hull}
\cppfile{geometry/convex.hull.cpp}
\subsection{Cut Polygon}
\cppfile{geometry/cut.polygon.cpp}
%\subsection{Bresenham}
%\cppfile{geometry/bresenham.cpp}
%\subsection{Rotate Matrix}
%\cppfile{geometry/rotate.cpp}
%\subsection{Interseccion de Circulos en n3log(n)}
%\cppfile{geometry/int.circs.cpp}
%
%\subsection{Cayley-Menger}
%{
%Permite calcular el volumen de un simplex (triángulo en k dimensiones) mediante el cálculo de una determinante.
%\cppfile[40-57]{geometry/cayley.cpp}
%}

\subsection{Heron's formula}
{
\( {\displaystyle A={\sqrt {s(s-a)(s-b)(s-c)}}} \),
\( {\displaystyle s={\frac {a+b+c}{2}}} \).
}

\section{DP Opt}
Observaciones:

$A[i][j]$ el menor $k$ que logra la solución óptima. En Knuth y D\&C la idea es aprovechar los rangos
determinados por este arreglo.

\subsection{Knuth}
{
    \textbf{Problema de ejemplo:} dado un palito de longitud $l$, con $n$ puntos en los que se puede cortar,
    determinar el costo mínimo para partir el palito en $n+1$ palitos unitarios (la DP se puede adaptar a $k$ 
    agregando un parámetro extra), donde hay un costo fijo por partir el rango $i, j$ que \emph{cumple la 
    condición suficiente}. Una función de costos que cumple es la distancia entre los extremos $j-i$. 
    El problema clásico de esta pinta es el del ABB óptimo.

    \textbf{Recurrencia original:} $dp[i][j] = min_{i < k < j} {dp[i][k] + dp[k][j]} + C[i][j]$ o bien 
		$dp[i][j] = min_{k < j}{dp[i - 1][k] + C[k][j]}$

    \textbf{Condición suficiente:} $A[i, j - 1] \leq  A[i, j] \leq A[i + 1, j]$

    Es decir, si saco un elemento a derecha el óptimo se mueve a izquierda o se mantiene, y 
    si saco un elemento a izquierda el óptimo se mueve a derecha o se mantiene.

    \textbf{Complejidad original:} $O(n^3)$

    \textbf{Complejidad optimizada:} $O(n^2)$

    \textbf{Solución:} iteramos por el tamaño $len$ del subarreglo (creciente), y para cada extremo izquierdo $l$,
    determinamos el extremo derecho $r=l+len$ e iteramos por los $k$ entre $A[l][r-1]$ y $A[l+1][r]$, actualizando la
    solución del estado actual.

    \cppfile{dp-opt/knuth.cpp}
}
\subsection{Divide \& Conquer}
{
    \textbf{Problema de ejemplo:} dado un arreglo de $n$ números con valores $a_1, a_1, \dots, a_n$, dividirlo
    en $k$ subarreglos, tal que la suma de los cuadrados del peso total de cada subarreglo es mínimo.

    \textbf{Recurrencia original:} $dp[i][j] = min_{k < j}{dp[i - 1][k] + C[k][j]}$ 

    \textbf{Condición suficiente:} $ A[i][j] \leq A[i][j + 1] $ o (normalmente más fácil de probar) 
    $ C[a][d] + C[b][c] \geq C[a][c] + C[b][d]$, con $a < b < c < d. $.

    La segunda condición suficiente es la intuición de que no conviene que los intervalos se contengan.

    \textbf{Complejidad original:} $O(kn^2)$

    \textbf{Complejidad optimizada:} $O(kn\log(n))$

    \textbf{Solución:} la idea es, para un $i$ determinado, partir el rango $[j_{left}, j_{right})$ al que pertenecen 
    los $j$ que queremos calcular a la mitad, determinar el óptimo y utilizarlo como límite para calcular los demás.
    Para implementar esto de forma sencilla, se suele utilizar la función recursiva $dp(i, j_{left}, j_{right}, opt_{left}, opt_{right})$
    que se encarga de, una vez fijado el punto medio $m$ del rango $[j_{left}, j_{right})$ iterar por los $k$ en $[j_{left}, j_{right})$ 
    para determinar el óptimo $opt$ para $m$, y continuar calculando $dp(i, j_{left}, m, opt_{left}, opt)$ y 
    $dp(i, m, j_{right}, opt, opt_{right})$.

    \cppfile{dp-opt/dc.cpp}
}

\section{Matemática}%%%%%%%%%%%%%%%%%%MATH%%%%%%%%%%%%%%%%%%%%%%%%%%%%%%%%
\subsection{Teoría de números}
{
\subsubsection{Funciones multiplicativas, función de Möbius}
{
Una funcion $f(n)$ es \textbf{multiplicativa} si para cada par de enteros coprimos $p$ y $q$ se cumple que \(f(pq) = f(p)f(q)\).

Si la función $f(n)$ es multiplicativa, puede evaluarse en un valor arbitrario conociendo los valores de la función en
sus factores primos: \(f(n) = f(p_1^{r_1})f(p_2^{r_2})\ldots f(p_k^{r_k})\).\\

La \textbf{función de Möbius} se define como:
\[\mu(n) = \begin{cases} 0 & d^2\ |\ n, \\ 1 & n=1,\\ (-1)^k & n = p_1p_2\cdots{p_k} .\end{cases}\]

\subsubsection{Teorema de Wilson}
{
\(\displaystyle (p-1)! \equiv -1 \;(\bmod\; p) \)
Siendo $p$ primo.
}
\subsubsection{Pequeño teorema de Fermat}
{
\(\displaystyle a^p \equiv a \;(\bmod\; p) \)
Siendo $p$ primo.
}
\subsubsection{Teorema de Euler}
{
\(\displaystyle a^{\varphi(n)} \equiv 1 \;(\bmod\; n) \)
}
}
\subsection{Combinatoria}
{
\subsubsection{Burnside's lemma}
{
Sea $G$ un grupo que actúa en un conjunto $X$. Para cada $g$ en $G$, sea $X^{g}$ el conjunto de elementos
en $X$ que son invariantes respecto a $g$, entonces el número de órbitas $|X/G|$ es:
\[ |X/G|={\frac {1}{|G|}}\sum _{{g\in G}}|X^{g}|.\]
Por ejemplo, si el grupo $G$ consiste de las operaciones de rotación, el conjunto $X$ son los posibles coloreos de un tablero, 
entonces el número de órbitas $|X/G|$ es el número de posibles coloreos de un tablero salvo rotaciones.
}

\subsubsection{Combinatorios}
{
\(\displaystyle \binom{n}{k} = \frac{n!}{k!(n-k)!} \)
\cppfile[12-23]{math/combinatorics.cpp}
}

\subsubsection{Lucas Theorem}
{
    \(\displaystyle {\binom {m}{n}}\equiv \prod_{i=0}^{k}{\binom {m{i}}{n_{i}}}{\pmod{p}}\)

where \({\displaystyle m=m_{k}p^{k}+m_{k-1}p^{k-1}+\cdots +m_{1}p+m_{0},}\)

and \({\displaystyle n=n_{k}p^{k}+n_{k-1}p^{k-1}+\cdots +n_{1}p+n_{0}}\)

\({\displaystyle {\tbinom {m}{n}}=0}\) if \(m < n\).


\cppfile[50-58]{math/combinatorics.cpp}
}
\subsubsection{Stirling}
{
\(\displaystyle \left\{{n \atop k}\right\} = \) cantidad de formas de particionar un conjunto de \(n\) elementos en \(m\) subconjuntos no vacíos.

\(\displaystyle \left\{{n+1\atop k}\right\} = k \left\{{ n \atop k }\right\} + \left\{{n\atop k-1}\right\}\)

for \(k > 0\) with initial conditions

\({\displaystyle \left\{{0 \atop 0}\right\}=1\quad {\mbox{ and }}\quad \left\{{n \atop 0}\right\}=\left\{{0 \atop n}\right\}=0}\)
for \(n > 0\).


\cppfile[33-40]{math/combinatorics.cpp}


\( {\displaystyle \left\{{n \atop k}\right\}={\frac {1}{k!}}\sum_{i=0}^{k}(-1)^{i}{\binom {k}{i}}(k-i)^{n}.} \)
}
\subsubsection{Bell}
{
\( {\displaystyle B_{n}=} \) cantidad de formas de particionar un conjunto de \(n\) elementos en subconjuntos no vacíos.

\( {\displaystyle B_0= B_1 = 1} \)

\( {\displaystyle B_{n+1}=\sum_{k=0}^{n} \binom{n}{k} B_k.} \) 

\( {\displaystyle B_{n}=\sum _{k=0}^{n}\left\{{n \atop k}\right\}.} \) 


\cppfile[42-48]{math/combinatorics.cpp}
} 
\subsubsection{Eulerian}
{
\( {\displaystyle A_{n, m}=} \) cantidad de permutaciones de 1 a n con m ascensos (m elementos mayores que el anterior).

\( {\displaystyle A(n,m)=(n-m)A(n-1,m-1)+(m+1)A(n-1,m).} \)
}
\subsubsection{Catalan}
{
\( {\displaystyle C_{n}=} \) cantidad de árboles binarios de n+1 hojas, en los que cada nodo tiene cero o dos hijos.

\({\displaystyle C_{n}={2n \choose n}-{2n \choose n-1}\quad {\mbox{ con }}n\geq 1.} \)

\( {\displaystyle C_{0}=1\quad {\mbox{y}}\quad C_{n+1}=\sum _{i=0}^{n}C_{i}\,C_{n-i}\quad {\mbox{con }}n\geq 0.} \)
}

}

%\subsection{Sumatorias conocidas}
%{
%$\sum_{i=0}^n\binom{n}{i}=2^n$
%
%$\sum_{i=0}^n i\binom{n}{i}=n*2^{n-1}$
%
%$\sum_{i=m}^n i = \frac{n(n+1)}{2} - \frac{m(m-1)}{2} = \frac{(n+1-m)(n+m)}{2}$
%
%$\sum_{i=0}^n i = \sum_{i=1}^n i = \frac{n(n+1)}{2}$
%
%$\sum_{i=0}^n i^2 = \frac{n(n+1)(2n+1)}{6} = \frac{n^3}{3} + \frac{n^2}{2} + \frac{n}{6}$
%
%$\sum_{i=0}^n i(i-1) = \frac{8}{6}(\frac{n}{2})(\frac{n}{2}+1)(n+1)$ (doubles) $\rightarrow$ Sino ver caso impar y par
%
%$\sum_{i=0}^n i^3 = \left(\frac{n(n+1)}{2}\right)^2 = \frac{n^4}{4} + \frac{n^3}{2} + \frac{n^2}{4} = \left[\sum_{i=1}^n i\right]^2$
%
%$\sum_{i=0}^n i^4 = \frac{n(n+1)(2n+1)(3n^2+3n-1)}{30} = \frac{n^5}{5} + \frac{n^4}{2} + \frac{n^3}{3} - \frac{n}{30}$
%
%$\sum_{i=0}^n i^p = \frac{(n+1)^{p+1}}{p+1} + \sum_{k=1}^p\frac{B_k}{p-k+1}{p\choose k}(n+1)^{p-k+1}$
%
%}
%
%\subsection{Ec. Característica}
%$a_0T(n)+a_1T(n-1)+...+a_kT(n-k)=0$
%
%$p(x)=a_0 x^k + a_1 x^{k-1} + ... + a_k$
%
%Sean $r_1,r_2,...,r_q$ las raíces distintas, de mult. $m_1, m_2, ..., m_q$
%
%$T(n)=\sum_{i=1}^q{\sum_{j=0}^{m_i - 1}c_{ij} n^j r_i^n}$

%Las constantes $c_{ij}$ se determinan por los casos base.
%\subsection{Aritmetica Modular}
%\cppfile{math/modular_arithmetic.cpp}
%\subsection{Exp. de Numeros Mod.}
%\cppfile{math/exp.mod.cpp}
%\subsection{Exp. de Matrices y Fibonacci en log(n)}
%\cppfile{math/exp.mat.cpp}
%\subsection{Primos}
%\cppfile{math/primes.cpp}
%\subsection{Factorizacion}
%Sea $n=\prod{p_i^{k_i}}$, fact(n) genera un map donde a cada $p_i$ le asocia su $k_i$
%\cppfile{math/factorize.cpp}
%\subsection{Divisores}
%\cppfile{math/divisors.cpp}
%\subsection{Euler's Phi}
%\cppfile[1-40]{math/phi_euler.cpp}
%\subsection{Phollard's Rho - Miller-Rabin}
%\cppfile[16-78]{math/phollards.rho.cpp}
%\subsection{GCD}
%\cppfile[1-2]{math/GCD-LCM.cpp}
%\subsection{LCM}
%\cppfile[3]{math/GCD-LCM.cpp}
\subsection{Euclides extendido}
Dados $a$ y $b$, encuentra $x$ e $y$ tales que $a*x + b*y = gcd(a, b)$.
\cppfile{math/extended.euclid.cpp}
\subsection{Inversos}
\cppfile{math/inversos.cpp}
\subsection{Ecuaciones diofánticas}
Basado en Euclides extendido. Dados $a$, $b$, y $r$ obtiene $x$ e $y$ tales que $a*x + b*y = r$, suponiendo que $gcd(a,b) | r$. Las soluciones son de la forma $(x, y) = (x_1 - b/gcd(a,b) * k_1, x_2 + a/gcd(a,b) * k_2)$ donde $x_1$ y $x_2$ son las soluciones particulares que obtuvo Euclides.
\cppfile{math/diophantine.cpp}
\subsection{Teorema Chino del Resto}
Dadas $k$ ecuaciones de la forma $a_i*x \equiv a_i \pmod {n_i}$, encuentra $x$ tal que es solución. Existe una única solución módulo $lcm(n_i)$.
\cppfile{math/chinese.cpp}
\subsection{Simpson}
\cppfile{math/simpson.cpp}
\subsection{Fraction}
\cppfile{math/frac.cpp}
%\subsection{Polinomio, Ruffini e interpolación de Lagrange}
%\textbf{Interpolación de Lagrange}: dados $n+1$ pares $(x_i, y_i)$  permite encontrar el polinomio de grado $n$ tal que $f(x_i) = y_i$. 
%
%\textbf{Explicación}: computa $P(x) = y_1*f_1(x) + y_2*f_2(x) + \ldots + y_{n+1}*f_{n+1}(x)$ donde $f_i(x) = \frac{g_i(x)}{g_i(x_i)}$, $g_i(x) = \frac{h(x)}{x - x_i}$ y $h(x) = (x-x_1)*(x-x_2)*\ldots*(x-x_{n+1})$.
%Usa Ruffini para la división de polinomios.
%
%Trucazo para computar en \(O(n)\): \(x_{i + 1} - x_i = x_{j + 1} - x_j\) para todo \(i, j < n\).
%
%\textbf{Ejemplo de problema}: tenés que calcular una respuesta que depende de un n y parece ser polinomial, conseguís un par de puntos e 
%intentás armar el polinomio (usando el algoritmo online u offline).
%
%\cppfile{math/polinomio.cpp}

\subsection{Matrices}
\cppfile{math/matrix.cpp}
\subsection{Determinante}
\cppfile{math/determinant.cpp}
\subsection{Sistemas de Ecuaciones Lineales - Gauss}
\cppfile{math/gauss.cpp}

\subsection{FFT y NTT}
\textbf{Base teórica (e intuición)}: \\
La \textbf{transformada de Fourier} mapea una función temporal a un dominio de frecuencias. \\
Podemos pensar que rotamos la función temporal alrededor de un círculo a diferentes frecuencias y
calculamos la magnitud del centro de masa de la figura resultante; la función del dominio de 
frecuencias representa este mapeo. \\
\[
    F_n = 
    \begin{bmatrix}
        1 & 1 & 1 & \ldots & 1 \\
        1 & \xi_n & \xi_n^2 & \dots & \xi_n^{n-1} \\
        1 & \xi_n^2 & \xi_n^4 & \dots & \xi_n^{2(n-1)} \\
        \vdots & \vdots & \vdots & \ddots & \vdots \\
        1 & \xi_n^{n-1} & \xi_n^{2(n-1)} & \dots & \xi_n^{(n-1)(n-1)}
    \end{bmatrix}
\]
\[
    y = F_nx
\]
Donde $\omega_n$ es una raíz primitiva n-ésima de la unidad y $\xi_n = w_n^{n-1}$.
La \textbf{transformada rápida de Fourier} se basa en que las raíces de la unidad cumplen la propiedad $ \omega_{2n}^{2} = \omega_{n} $. Por lo tanto:
\[
    F_n = 
    \begin{bmatrix}
        F_{n/2} & D_{n/2}F_{n/2} \\
        F_{n/2} & -D_{n/2}F_{n/2}
    \end{bmatrix}
    P_n 
\]
donde: 
\[
    D_{n/2} = 
    \begin{bmatrix}
        1 \\
        & \xi_n \\
        & & \xi_n^2 \\
        & & & \ddots \\
        & & & & \xi_n^{n/2-1}
    \end{bmatrix}
\]
    y
\[
    P_n^T = 
    \begin{bmatrix}
        e_0 & e_2 & e_4 & \ldots & e_{n-2} & e_1 & e_3 & \ldots & e_{n-1}
    \end{bmatrix}
\]  
\textbf{NTT}: es un algoritmo más lento pero más preciso para calcular la DFT,
ya que trabaja con enteros módulo un primo $m$. \\
El módulo $m$ debe ser un primo de la forma $m = c2^k + 1$.
Para encontrar la raíz $2^k$-ésima de la unidad $r$: $r = g ^ c$, donde $g$ es una raíz primitiva de $p$ (número tal que si lo elevamos a diferentes potencias recorremos todos los demás). \\
Valores tradicionales: $m = 998244353$ y $r = 3$, $m = 2305843009255636993$ y $r = 5$ (este último da overflow, se podría fixear). \\

\textbf{Operaciones}: \\
Es mucho más fácil realizar ciertas operaciones en un dominio de frecuencias:
\begin{itemize}
	\item Multiplicar en $O(n\log(n))$: simplemente multiplicar punto a punto.

    \item Invertir en $O(n\log(n))$: asumiendo $B(0) \neq 0$, existe una serie infinita $C(x)$ que es inverso del polinomio.
    Aprovechando ciertas propiedades del producto $B(x)C(x)$ ($b_0c_0 = 1$ y el resto de los coeficientes resultantes son $0$),
    podemos ir despejando el inverso. Es posible aplicar Divide and Conquer notando la relación entre los primeros $n/2$ términos
    del inverso y los siguientes $n/2$.

	\item Dividir en $O(n\log(n))$: resulta más fácil dividir los polinomios reversos (ya que un polinomio y su reverso son casi iguales,
	y no hace falta considerar resto de la división de los reversos).
	
    \item Multievaluar en $O(n\log^2(n))$: evaluar un polinomio $A(x)$ en $x_1$ es lo mismo que dividir $A(x)$ por $x - x_1$ y evaluar el resto
    $R(x)$ en $x_1$. Para múltiples puntos, podemos utilizar una estrategia estilo Divide and Conquer.

    \item Interpolar en $O(n\log^2(n))$: para interpolar se utilizan los polinomios de Lagrange (ver interpolación de Lagrange, 
    $A(x)=\sum_{i=1}^{n} y_i\frac{1}{p_i(x_i)}p_i(x)$ y $p_i(x) = \frac{p(x)}{x-x_i}$). Para poder computarlos rápidamente, aprovechamos que
    $p'(x_i)=p_i(x_i)$ (podemos computar la derivada y evaluar con multievaluación) y utilizamos una estrategia estilo Segment Tree para
    generar los polinomios rápidamente (notando que si mantenemos los polinomios para dos conjuntos de puntos es fácil unirlos).

\end{itemize}

\cppfile{math/fft_ntt.cpp}
\cppfile{math/fft_ntt_operations.cpp}

\subsection{Programación lineal: Simplex}
\textbf{Introducción} \\
Permite maximizar cierta función lineal dado un conjunto de restricciones lineales.
\textbf{Algoritmo} \\
El algoritmo opera con programas lineales en la siguiente forma canónica: maximizar \(z = c^{T}x \) sujeta a \(Ax \leq b, x \geq 0\).  \\
Por ejemplo, si \(c=(2, -1)\), \(A=\begin{bmatrix} 1 & 0 \end{bmatrix}\) y \(b=(5)\), 
buscamos maximizar \(z = 2x_1 - x_2\) sujeta a \(x_1 \leq 5\) y \(x_i \geq 0\). \\
\textbf{Detalles implementativos} \\
Canonizar si hace falta.\\
Para obtener soluciones negativas, realizar el cambio de variable \(x_i = x_i' + \text{INF}\).\\
Si la desigualdad no incluye igual, solo menor, \textbf{no usar epsilon} al agregarla. Esto ya es considerado por el código.
\cppfile{math/simplex.cpp}

%\subsection{Tablas y cotas (Primos, Divisores, Factoriales, etc)}
%%\subsubsection{
%\paragraph{Factoriales} \ \\
%\begin{tabular}{l|l}
%0! =	1             & 11! = 39.916.800  \\
%1! =	1             & 12! =	479.001.600	($\in \mathtt{int}$)\\
%2! =	2             & 13! =	6.227.020.800	\\
%3! =	6             & 14! =	87.178.291.200	\\
%4! =	24            & 15! =	1.307.674.368.000	\\
%5! =	120   			  & 16! =	20.922.789.888.000	\\
%6! =	720           & 17! =	355.687.428.096.000	\\
%7! =	5.040	        & 18! =	6.402.373.705.728.000	\\
%8! =	40.320	      & 19! =	121.645.100.408.832.000	\\
%9! =	362.880       & 20! =	2.432.902.008.176.640.000	($\in \mathtt{tint}$) \\
%10! =	3.628.800     & 21! =	51.090.942.171.709.400.000
%\end{tabular}
% 
%max signed tint = 9.223.372.036.854.775.807 \\
%max unsigned tint = 18.446.744.073.709.551.615
%%\subsubsection{
%\paragraph{Primos} \ \\
%2 3 5 7 11 13 17 19 23 29
%31 37 41 43 47 53 59 61 67 71
%73 79 83 89 97 101 103 107 109 113
%127 131 137 139 149 151 157 163 167 173
%179 181 191 193 197 199 211 223 227 229
%233 239 241 251 257 263 269 271 277 281
%283 293 307 311 313 317 331 337 347 349
%353 359 367 373 379 383 389 397 401 409
%419 421 431 433 439 443 449 457 461 463
%467 479 487 491 499 503 509 521 523 541
%547 557 563 569 571 577 587 593 599 601
%607 613 617 619 631 641 643 647 653 659
%661 673 677 683 691 701 709 719 727 733
%739 743 751 757 761 769 773 787 797 809
%811 821 823 827 829 839 853 857 859 863
%877 881 883 887 907 911 919 929 937 941
%947 953 967 971 977 983 991 997 1009 1013
%1019 1021 1031 1033 1039 1049 1051 1061 1063 1069
%1087 1091 1093 1097 1103 1109 1117 1123 1129 1151
%1153 1163 1171 1181 1187 1193 1201 1213 1217 1223
%1229 1231 1237 1249 1259 1277 1279 1283 1289 1291
%1297 1301 1303 1307 1319 1321 1327 1361 1367 1373
%1381 1399 1409 1423 1427 1429 1433 1439 1447 1451
%1453 1459 1471 1481 1483 1487 1489 1493 1499 1511
%1523 1531 1543 1549 1553 1559 1567 1571 1579 1583
%1597 1601 1607 1609 1613 1619 1621 1627 1637 1657
%1663 1667 1669 1693 1697 1699 1709 1721 1723 1733
%1741 1747 1753 1759 1777 1783 1787 1789 1801 1811
%1823 1831 1847 1861 1867 1871 1873 1877 1879 1889
%1901 1907 1913 1931 1933 1949 1951 1973 1979 1987
%1993 1997 1999 2003 2011 2017 2027 2029 2039 2053
%2063 2069 2081
 
%2083 2087 2089 2099 2111 2113 2129
%2131 2137 2141 2143 2153 2161 2179 2203 2207 2213
%2221 2237 2239 2243 2251 2267 2269 2273 2281 2287
%2293 2297 2309 2311 2333 2339 2341 2347 2351 2357
%2371 2377 2381 2383 2389 2393 2399 2411 2417 2423
%2437 2441 2447 2459 2467 2473 2477 2503 2521 2531
%2539 2543 2549 2551 2557 2579 2591 2593 2609 2617
%2621 2633 2647 2657 2659 2663 2671 2677 2683 2687
%2689 2693 2699 2707 2711 2713 2719 2729 2731 2741
%2749 2753 2767 2777 2789 2791 2797 2801 2803 2819
%2833 2837 2843 2851 2857 2861 2879 2887 2897 2903
%2909 2917 2927 2939 2953 2957 2963 2969 2971 2999
%3001 3011 3019 3023 3037 3041 3049 3061 3067 3079
%3083 3089 3109 3119 3121 3137 3163 3167 3169 3181
%3187 3191 3203 3209 3217 3221 3229 3251 3253 3257
%3259 3271 3299 3301 3307 3313 3319 3323 3329 3331
%3343 3347 3359 3361 3371 3373 3389 3391 3407 3413
%3433 3449 3457 3461 3463 3467 3469 3491 3499 3511
%3517 3527 3529 3533 3539 3541 3547 3557 3559 3571\\
%\paragraph{Primos cercanos a $10^n$}\ \\
%9941 9949 9967 9973 10007 10009 10037 10039 10061 10067 10069 10079\\
%99961 99971 99989 99991 100003 100019 100043 100049 100057 100069\\
%999959 999961 999979 999983 1000003 1000033 1000037 1000039\\
%9999943 9999971 9999973 9999991 10000019 10000079 10000103 10000121\\
%99999941 99999959 99999971 99999989 100000007 100000037 100000039 100000049\\
%999999893 999999929 999999937 1000000007 1000000009 1000000021 1000000033
% 
%\paragraph{Cantidad de primos menores que $10^n$}\ \\
%$\pi(10^1)$ = 4 ;
%$\pi(10^2)$ = 25 ;
%$\pi(10^3)$ = 168 ;
%$\pi(10^4)$ = 1229 ;
%$\pi(10^5)$ = 9592 \\
%$\pi(10^6)$ = 78.498 ;
%$\pi(10^7)$ = 664.579 ;
%$\pi(10^8)$ = 5.761.455 ;
%$\pi(10^9)$ = 50.847.534 \\
%$\pi(10^{10})$ = 455.052,511 ;
%$\pi(10^{11})$ = 4.118.054.813 ;
%$\pi(10^{12})$ = 37.607.912.018 \\
%\textbf{Observación:} Una buena aproximación es $x/ln(x)$.
%
% Fuente: http://primes.utm.edu/howmany.shtml#table
%
%
%\subsubsection{Divisores}
%\paragraph{Divisores} \ \\
%Cantidad de divisores ($\sigma_0$) para \emph{algunos} $n / \neg\exists n'<n, \sigma_0(n') \geqslant \sigma_0(n)$ \\
%\textbf{Referencias:} $\sigma_0(10^9) = 1344$ y $\sigma_0(10^{18}) = 103680$ \\
%$\sigma_0(60)$ = 12 ; $\sigma_0(120)$ = 16 ; $\sigma_0(180)$ = 18 ; $\sigma_0(240)$ = 20 ; $\sigma_0(360)$ = 24 \\
%$\sigma_0(720)$ = 30 ; $\sigma_0(840)$ = 32 ; $\sigma_0(1260)$ = 36 ; $\sigma_0(1680)$ = 40 ; $\sigma_0(10080)$ = 72 \\ $\sigma_0(15120)$ = 80 ; $\sigma_0(50400)$ = 108 ; $\sigma_0(83160)$ = 128 ; $\sigma_0(110880)$ = 144 \\
%$\sigma_0(498960)$ = 200 ; $\sigma_0(554400)$ = 216 ; $\sigma_0(1081080)$ = 256 ; $\sigma_0(1441440)$ = 288  $\sigma_0(4324320)$ = 384 ; $\sigma_0(8648640)$ = 448 \\
%\textbf{Observación:} Una buena aproximación es $x^{1/3}$. \\ \\
%%
%Suma de divisores ($\sigma_1$) para \emph{algunos} $n / \neg\exists n'<n, \sigma_1(n') \geqslant \sigma_1(n)$ \\
%$\sigma_1(96)$ = 252 ; $\sigma_1(108)$ = 280 ; $\sigma_1(120)$ = 360 ; $\sigma_1(144)$ = 403 ; $\sigma_1(168)$ = 480 \\
%$\sigma_1(960)$ = 3048 ; $\sigma_1(1008)$ = 3224 ; $\sigma_1(1080)$ = 3600 ; $\sigma_1(1200)$ = 3844 \\
%$\sigma_1(4620)$ = 16128 ; $\sigma_1(4680)$ = 16380 ; $\sigma_1(5040)$ = 19344 ; $\sigma_1(5760)$ = 19890 \\
%$\sigma_1(8820)$ = 31122 ; $\sigma_1(9240)$ = 34560 ; $\sigma_1(10080)$ = 39312 ; $\sigma_1(10920)$ = 40320 \\
%$\sigma_1(32760)$ = 131040 ; $\sigma_1(35280)$ = 137826 ; $\sigma_1(36960)$ = 145152 ; $\sigma_1(37800)$ = 148800 \\
%$\sigma_1(60480)$ = 243840 ; $\sigma_1(64680)$ = 246240 ; $\sigma_1(65520)$ = 270816 ; $\sigma_1(70560)$ = 280098 \\
%$\sigma_1(95760)$ = 386880 ; $\sigma_1(98280)$ = 403200 ; $\sigma_1(100800)$ = 409448  \\
%$\sigma_1(491400)$ = 2083200 ; $\sigma_1(498960)$ = 2160576 ; $\sigma_1(514080)$ = 2177280 \\
%$\sigma_1(982800)$ = 4305280 ; $\sigma_1(997920)$ = 4390848 ; $\sigma_1(1048320)$ = 4464096 \\
%$\sigma_1(4979520)$ = 22189440 ; $\sigma_1(4989600)$ = 22686048 ; $\sigma_1(5045040)$ = 23154768 \\
%$\sigma_1(9896040)$ = 44323200 ; $\sigma_1(9959040)$ = 44553600 ; $\sigma_1(9979200)$ = 45732192
%
%


\section{Grafos}%%%%%%%%%%%%%%%%%%GRAFOS%%%%%%%%%%%%%%%%%%%%%%%%%%%%
\subsection{Teoremas y fórmulas}
{
\subsubsection{Teorema de Pick}
{
$A=I+\frac{B}{2}-1$

Donde \(A\) es el área, \(I\) es la cantidad de puntos interiores, y \(B\) la cantidad de puntos en el borde.
}

\subsubsection{Formula de Euler}
{
$v-e+f=k+1$

Donde \(v\) es la cantidad de vértices, \(e\) la cantidad de arcos, \(f\) la cantidad de caras y \(k\) la cantidad de componentes conexas.
}
}
%\subsection{Dijkstra}
%\cppfile{graphs/dijkstra.cpp}
\subsection{Bellman-Ford}
\cppfile{graphs/bellman.ford.cpp}
%\subsection{Floyd-Warshall}
%\cppfile{graphs/floyd.warshall.cpp}
\subsection{Kruskal}
\cppfile{graphs/kruskal.cpp}
%\subsection{Prim}
%\cppfile[23-40]{graphs/prim.cpp}
\subsection{2-SAT + Tarjan SCC}
\cppfile{graphs/2sat.cpp}
%\subsection{Tarjan}
%\cppfile{graphs/tarjan.cpp}
%\subsection{Kosaraju}
%\cppfile[23-72]{graphs/kosaraju.cpp}
\subsection{Articulation Points}
\cppfile{graphs/articulaciones.cpp}
\subsection{Comp. Biconexas y Puentes}
\cppfile{graphs/biconexas.bridge.cpp}
\subsection{LCA + Climb}
\cppfile{graphs/lca.climb.cpp}
%\subsection{Union Find} 
%\cppfile{graphs/union.find.cpp}

%\subsection{Splay Tree + Link-Cut Tree}
%\textbf{Definición:} Splay Tree es un \emph{binary search tree} eficiente.
%
%\textbf{Operaciones:} soporta operaciones en $O(\log{}n)$ amortizado.
%\begin{itemize}
%    \item $splay(X)$: lleva el nodo $X$ a la raíz del árbol (manteniendo el orden de los nodos). 
%    Para esto, se realizan rotaciones de forma tal que el árbol "quede más balanceado", resultando
%    en una complejidad amortizada de $O(\log{}n)$.
%    \item $search(X)$: igual que en cualquier árbol binario de búsqueda, pero después se ejecuta $splay(X)$.
%    \item $split(X)$ y $merge(X, Y)$: usando $search$, funcionan igual que en Treap.
%    \item $insert(X)$ y $erase(X)$: usando $split$ y $merge$, igual que en Treap.
%    \item $reverse()$ y otras operaciones asociativas: usando lazy-propagation, similar a Treap.
%    \item Mantener una lista en vez de un conjunto: igual que en Treap.
%\end{itemize}
%
%\textbf{Definición:} Link-Cut Tree es una estructura que mantiene un \emph{bosque de árboles con raíz}.
%
%\textbf{Conceptos base:}
%\begin{itemize}
%    \item Hijo preferido de $X$: es la raíz del subárbol en el cual se realizó el último acceso 
%    entre los descendientes de $X$, o ninguno si el último acceso fue en $X$.
%    \item Camino preferido: se representan con un \emph{splay tree}, además se almacena el padre 
%    del nodo más alto del camino (manteniendo así toda la info del árbol original)
%    \item $access(X)$: es la operación básica de la estructura, se hace cada vez que realizamos una
%    operación sobre el nodo $X$. Se encarga de mantener actualizado al "hijo preferido" de cada ancestro
%    de $X$, realizando operaciones en los splay trees correspondientes. Primero hace $splay(X)$ y desconecta
%    a su hijo derecho. Luego, iterativamente: hace $splay(Y)$ con $Y$ padre del camino de $X$, corta 
%    al hijo derecho de $Y$, luego concatena el camino preferido de $Y$ con el de $X$, y finalmente rota 
%    a $X$ y continúa.
%\end{itemize}
%
%\textbf{Operaciones:} soporta operaciones en $O(\log{}n)$ amortizado.
%\begin{itemize}
%    \item $link(X, Y)$: hacer a $X$ hijo de $Y$.
%    \item $cut(X)$: desconectar a $X$ de su padre.
%    \item $makeRoot(X)$: hace a $X$ raíz de su árbol.
%    \item $getRoot(X)$: devuelve la raíz del árbol de $X$.
%    \item $lca(X,Y)$: LCA tradicional.
%    \item $lift(X,k)$: k-ésimo ancestro.
%    \item $update(X,v)$: cambiar a $v$ el valor asociado a $X$.
%    \item $aggregate(X,Y)$: resultado de una operación asociativa en los nodos del camino de $X$ a $Y$.
%\end{itemize}
%
%\cppfile{graphs/link.cut.tree.splay.tree.cpp}
%
%\subsection{Heavy Light Decomposition}
%\cppfile{graphs/heavy.light.cpp}
%\subsection{Centroid Decomposition}
%\cppfile[3-20]{graphs/centroid.cpp}
%\subsection{Euler Cycle}
%\cppfile{graphs/euler.cpp}
%\subsection{Diametro árbol}
%\cppfile{graphs/diametro.cpp}
%\subsection{Chu-liu}
%\cppfile{graphs/chuliu.cpp}
%\subsection{Dynamic Conectivity}
%\textbf{Definición:} permite realizar queries sobre un grafo dinámico al que se le pueden agregar y quitar aristas.
%
%\textbf{Explicación: } procesa las queries (y los updates) offline, con una estrategia muy similar a la de la búsqueda
%binaria en paralelo: pensar que los arcos están presentes en cierto intervalo de tiempo, y que solo incluimos los arcos
%que contienen totalmente al intervalo que estamos considerando (a medida que se mueven los extremos). Al igual
%que en la búsqueda binaria en paralelo, se puede ver que se forma un árbol binario en el que se realiza una cantidad
%de operaciones lineal en cada nivel.
%\cppfile{graphs/dynamic.conectivity.cpp}


\section{Flujo}%%%%%%%%%%%%%%%%%%FLOW%%%%%%%%%%%%%%%%%%%%%%%%%%%%
\subsection{Trucazos generales}
\begin{itemize}
\item \textbf{Corte mínimo}: aquellos nodos alcanzables desde $S$ forman un conjunto, los demás forman el otro conjunto.
En Dinic's: vertices con $dist[v] >= 0$ (del lado de $S$) vs.  $dist[v] == -1$ (del lado del $T$).
\item \textbf{Para grafos bipartitos}: sean $V_1$ y $V_2$ los conjuntos más próximos a $S$ y a $T$ respectivamente.
	\begin{itemize}
	\item \textbf{Matching}: para todo $v_1 \in V_1$ tomar las aristas a vértices en $V_2$ con flujo positivo ($edge.f > 0$).
	\item \textbf{Min. Vertex Cover}: unión de vértices $v_1 \in V_1$ tales que son inalcanzables ($dist[v_1] == -1$),
	y vértices $v_2 \in V_2$ tales que son alcanzables ($dist[v_2] > 0$).
	\item \textbf{Max. Independent Set}: tomar vértices no tomados por el Min. Vertex Cover.
	\item \textbf{Max. Clique}: construir la red $G'$ (red complemento) y encontrar Max. Independent Set.
	\item \textbf{Min. Edge Cover}: tomar las aristas del Matching y para todo vértice no cubierto hasta el momento, 
	tomar cualquier arista incidente.
    \item \textbf{Konig's theorem}: $|$minimum vertex cover$|$ = $|$maximum matching$| \Leftrightarrow |$maximum independent set$| + |$maximum matching$| = |$vertices$|$.\\
	\end{itemize}
\end{itemize}

%\subsection{Ford Fulkerson}
%\textbf{Complejidad: } $O(fE)$.
%\textbf{Algoritmo: } cambiar BFS por DFS en Edmonds Karp.
%
%\subsection{Edmonds Karp}
%\textbf{Complejidad: } $O(VE^2)$.
%\cppfile{flow/edmonds.karp.cpp}

\subsection{Dinic}
\textbf{Complejidad: } $O(V^2E)$ en general. $O(min(E^{3/2}, V^{2/3}E))$ con capacidades unitarias. $O(\sqrt{V}E)$ en matching bipartito (se lo llama Hopcroft–Karp algorithm) y en cualquier otra red unitaria (indegree = outdegree = 1 para cada vértice excepto S y T). 
\cppfile{flow/dinic.cpp}

%\subsection{Maximum matching}
%\textbf{Complejidad}: $O(VE)$.
%\cppfile{flow/matching.cpp}
%
%\subsection{Hungarian}
%\textbf{Complejidad: } $O(V^3)$.
%\cppfile{flow/hungarian.cpp}

\subsection{Min-cost Max-flow}
\textbf{Algoritmo}: tira camino mínimo hasta encontrar el flujo buscado. Usa SPFA (Bellman-Ford más inteligente, con mejor
tiempo promedio) porque resulta en la mejor complejidad.

\textbf{Complejidad}: $O(V^2E^2)$.
\cppfile{flow/min.cost.max.flow.cpp}

\subsection{Flujo con demandas}
\textbf{Problema}: se pide que $d(e) \leq f(e) \leq c(e)$. \\
\textbf{Flujo arbitrario}: transformar red de la siguiente forma. Agregar nueva fuente $s'$ y nuevo sumidero $t'$,
arcos nuevos de $s'$ a todos los demás nodos, arcos nuevos desde todos los nodos a $t'$, y un arco de $t$ a $s$.
Definimos la nueva función de capacidad $c'$ como:
\begin{itemize}
\item $c'((s',v)) = \sum\nolimits_{u \in V} d((u,v))$ para cada arco $(s',v)$.
\item $c'((v,t')) = \sum\nolimits_{w \in V} d((v,w))$ para cada arco $(v,t')$.
\item $c'((u,v)) = c((u,v)) - d((u,v))$ para cada arco $(u,v)$ en la red original.
\item $c'((t,s)) = \infty$
\end{itemize}
\textbf{Flujo mínimo}: hacer búsqueda binaria sobre la capacidad de la arco $(t, s)$, viendo que se satisfaga la demanda.

\section{Template}%%%%%%%%%%%%%%%%%%TEMPLATE%%%%%%%%%%%%%%%%
\cppfile{others/template.cpp}

\section{vimrc}
\cppfile{others/vimrc-short}

\section{Misc}
\cppfile{others/truquillos.cpp}

%\subsection{Fast read}
%\cppfile{others/fast.read.cpp}

%\section{Ayudamemoria}%%%%%%%%%%%%%%%%%%AYUDAMEMORIA%%%%%%%%%%%%%%%%
%\subsection*{Cant. decimales} 	
%\begin{code}
%#include <iomanip>
%cout << setprecision(2) << fixed;
%\end{code}
%\subsection*{Rellenar con espacios(para justificar)}
%\begin{code}
%#include <iomanip>
%cout << setfill(' ') << setw(3) << 2 << endl;
%\end{code}
%\subsection*{Comparación de Doubles}
%\begin{code}
%const double EPS = 1e-9;
%x == y	<=> fabs(x-y) < EPS
%x >  y	<=> x > y + EPS
%x >= y	<=> x > y - EPS
%\end{code}
%\subsection*{Limites}
%\begin{code}
%#include <limits>
%numeric_limits<T>
%	::max()
%	::min()
%	::epsilon()
%\end{code}
%\subsection*{Muahaha}
%\begin{code}
%#include <signal.h>
%void divzero(int p){
%	while(true);}
%void segm(int p){
%	exit(0);}
%//in main
%signal(SIGFPE, divzero);
%signal(SIGSEGV, segm);
%\end{code}
%\subsection*{Mejorar velocidad 2}
%\begin{code}
%//Solo para enteros positivos
%inline void Scanf(int& a){
%	char c = 0;
%	while(c<33) c = getc(stdin);
%	a = 0;
%	while(c>33)	a = a*10 + c - '0', c = getc(stdin);
%}
%\end{code}
%\subsection*{Leer del teclado}
%\begin{code}
%freopen("/dev/tty", "a", stdin);
%\end{code}
%\subsection*{Iterar subconjunto}
%\begin{code}
%for(int sbm=bm; sbm; sbm=(sbm-1)&bm)
%\end{code}
%\subsection*{File setup}
%\begin{code}
%// tambien se pueden usar comas: {a, x, m, l}
%touch {a..l}.in; tee {a..l}.cpp < template.cpp
%\end{code}
%\subsection*{Releer String}
%\begin{code}
%string s; int n;
%getline(cin, s);
%stringstream leer(s);
%while(leer >> n){
%	// do something ...
%}
%\end{code}
\end{document}
